\section*{अध्याय \ref{ch:b1:clastic-sediments} --- खंडजात तलछट और उनके परिवेश}

\begin{solution}{exo:b1:clastic-sediments:1}
$\varphi = -\log_2 d$: \qty{0.18}{mm} से $2.47$ आता है, एक बारीक बालू (2 से 3);
\qty{0.75}{mm} से $0.42$, एक मोटी बालू (0 से 1); \qty{3}{mm} से
$-1.58$, एक कणिका ($-1$ से $-2$); \qty{0.012}{mm} से $6.38$, एक बारीक गाद (6
से 7); \qty{150}{mm} से $-7.23$, एक रोड़ी ($-6$ से $-8$)।
\end{solution}

\begin{solution}{exo:b1:clastic-sediments:2}
$M_z = (1.05 + 1.71 + 2.43)/3 = 1.73\,\varphi$, अर्थात् $2^{-1.73} =
\qty{0.30}{mm}$, एक मध्यम बालू। $\sigma_I = (2.43 - 1.05)/4 + (2.98 - 0.62)/6.6
= 0.345 + 0.358 = 0.70\,\varphi$: मध्यम रूप से अच्छी तरह छँटी (0.50 से 0.71), मध्यम रूप से छँटी की सीमा से ठीक नीचे।
\end{solution}

\begin{solution}{exo:b1:clastic-sediments:3}
$q_b = (1 - 0.4) \times 0.02 \times (\num{0.01}/60)/2 = \qty{1.0e-6}{m^2.s^{-1}}$
ठोस कणों का। एक दिन में, $\num{1.0e-6} \times 86400 = \qty{0.0864}{m^3}$ प्रति
मीटर शिखर, या $0.0864 \times 2650 = \qty{229}{kg}$ क्वार्ट्ज।
\end{solution}

\begin{solution}{exo:b1:clastic-sediments:4}
(क)~50 नॉट $50 \times 1852/3600 = \qty{25.7}{m.s^{-1}}$ हैं; 12 नॉट,
\qty{6.17}{m.s^{-1}}। (ख)~$\qty{100000}{m}/\qty{6.17}{m.s^{-1}} =
\qty{16200}{s}$, लगभग \qty{4.5}{h}। (ग)~480 मील \qty{772}{km} हैं, 400 मील
\qty{644}{km}।
\end{solution}

\begin{solution}{exo:b1:clastic-sediments:5}
(क)~$(\varphi_{84} - \varphi_{16})/2 = 2 \times 0.9945\,\sigma/2 =
0.9945\,\sigma$, और $(\varphi_{95} - \varphi_5)/3.3 = 2 \times
1.6449\,\sigma/3.3 = 0.9969\,\sigma$: दोनों $\sigma$ का अनुमान एक-तिहाई प्रतिशत के भीतर लगाते हैं। (ख)~एक वास्तविक तलछट फ़ाई इकाइयों में सामान्य होना ज़रूरी नहीं: बहुत मोटे या बहुत बारीक कणों का एक छोटा मिश्रण $\varphi_5$ या
$\varphi_{95}$ को बहुत और $\varphi_{16}$ तथा $\varphi_{84}$ को थोड़ा बदलता है, इसलिए दूसरा अनुमान बढ़ता है जबकि पहला नहीं। दोनों मिश्रणों के लिए सबसे अधिक भिन्न होते हैं, जैसे कुछ बजरी और कीचड़ वाली बालू, या द्विबहुलक निक्षेप।
\end{solution}

\begin{solution}{exo:b1:clastic-sediments:6}
(क)~$\sin\theta = T/L = 1.5/15 = 0.1$, इसलिए $\theta = \ang{5.74}$। (ख)~$a =
c\tan\theta = 0.5 \times 0.1005 = \qty{0.050}{cm}$ प्रति मिनट, या
\qty{3.0}{cm} प्रति घंटा। (ग)~सेंटीमीटर प्रति घंटे पर उठता तल निलंबन से गिरती बालू बहुत तेज़ी से पाता है: एक भारी लदी धारा जो धीमी हो रही है, जैसे घटती बाढ़ या एक \omterm{def:b1:clastic-sediments:turbidite}{गंदलाव-धारा} की पूँछ।
\end{solution}

\begin{solution}{exo:b1:clastic-sediments:7}
(क)~$(350 + 20 + 5 + 340 + 15 + 30)/6 = 760/6 = 126.7$: एक दक्षिण-पूर्वी
दिशा, यद्यपि हर पाठन उत्तर से \ang{20} के भीतर है। (ख)~$S =
\sum\sin\alpha_i = 0.672$, $C = \sum\cos\alpha_i = 5.692$, इसलिए $\bar\alpha =
\operatorname{atan2}(0.672, 5.692) = \ang{6.7}$, और $\bar R =
\sqrt{0.672^2 + 5.692^2}/6 = 5.732/6 = 0.955$। (ग)~अंकगणितीय माध्य 350 और 005 को दूर मानता है, जबकि वृत्त पर वे \ang{15} दूर हैं: दिगंशों का माध्य इकाई सदिशों के रूप में लेना चाहिए। बड़ा $\bar R$ दिखाता है कि पाठन \ang{7} के चारों ओर कसकर समूहित हैं।
\end{solution}

\begin{solution}{exo:b1:clastic-sediments:8}
A में \qty{5}{\percent} से अधिक मृत्तिका है: अपरिपक्व, और संघटनात्मक रूप से भी अपरिपक्व, अपने फ़ेल्डस्पार और चट्टान-टुकड़ों के साथ; अपने स्रोत के निकट शीघ्र जमा, एक जलोढ़ पंखे या बाढ़-मैदान पर, या एक \omterm{def:b1:clastic-sediments:turbidite}{गंदलाव-धारा} से। B में \qty{5}{\percent} से कम मृत्तिका है और अच्छी तरह छँटी है ($\sigma_I < 0.5\,\varphi$), पर उसके कण उप-कोणीय हैं ($\rho < 3$): परिपक्व, उदाहरण के लिए एक नदी-पट्टी या धाराओं से फटकी शेल्फ़ की बालू। C में मृत्तिका नहीं है, बहुत अच्छी तरह छँटी है और उसके कण गोल हैं ($\rho > 3$): अति-परिपक्व, और \qty{98}{\percent} क्वार्ट्ज के साथ संघटनात्मक रूप से भी परिपक्व, जैसे एक तट पर या मरुस्थली टीला-क्षेत्र में।
\end{solution}

\begin{solution}{exo:b1:clastic-sediments:9}
(क)~$1.5/\sqrt{9.81 \times 0.3} = 0.87$; $1.0/\sqrt{9.81 \times 2} = 0.23$।
(ख)~$h = U^2/g = 1.5^2/9.81 = \qty{0.23}{m}$; $U = \sqrt{9.81 \times 2} =
\qty{4.43}{m.s^{-1}}$। (ग)~दिए गए वेग पर फ़ूड संख्या गहराई घटने के साथ बढ़ती है: पतली, तेज़ जल-चादरें, जैसे तट पर एक लहर का पश्च-प्रवाह, आसानी से 1 के निकट फ़ूड संख्याएँ पाती हैं, जबकि कई मीटर गहरी नदी को कई मीटर प्रति सेकंड चाहिए, जो बाढ़ों के अतिरिक्त कुछ ही नदियाँ पहुँचती हैं।
\end{solution}

\begin{solution}{exo:b1:clastic-sediments:10}
(क)~मध्य-बिंदुओं $-0.5, 0.5, \ldots, 4.5$ के साथ: $m = (-1 + 4 + 45 + 100 + 52.5
+ 22.5)/100 = 2.23\,\varphi$; प्रतिशतों से भारित वर्ग-विचलनों का योग \num{107.71} है, इसलिए $s = \sqrt{1.0771} = 1.04\,\varphi$:
खराब छँटी। (ख)~संचयी मोटा प्रतिशत $-1, 0, \ldots, 5\,\varphi$ पर 0, 2, 10, 40, 80, 95 और
100 है। तब $\varphi_5 = 0 + 3/8 = 0.375$,
$\varphi_{16} = 1 + 6/30 = 1.20$, $\varphi_{50} = 2 + 10/40 = 2.25$,
$\varphi_{84} = 3 + 4/15 = 3.27$, $\varphi_{95} = 4.0$। $M_z = (1.20 + 2.25 +
3.27)/3 = 2.24\,\varphi$ और $\sigma_I = (3.27 - 1.20)/4 + (4.0 - 0.375)/6.6 =
0.517 + 0.549 = 1.07\,\varphi$। (ग)~दोनों विधियाँ माध्य पर लगभग \num{0.01} और \omterm{def:b1:clastic-sediments:sorting}{छँटाव} पर \num{0.03} तक सहमत हैं: एक बारीक बालू
($2.24\,\varphi$ \qty{0.21}{mm} है), खराब छँटी।
\end{solution}

\begin{solution}{exo:b1:clastic-sediments:11}
(क)~$U = \sqrt{0.02 \times 9.81 \times 50 \times 0.005/0.003} =
\qty{4.04}{m.s^{-1}}$, या 7.9 नॉट। (ख)~$(\Delta\rho/\rho)\,h = U^2
C_f/(gS) = 25.7^2 \times 0.003/(9.81 \times 0.005) = \qty{40}{m}$: उदाहरण के लिए
\qty{400}{m} मोटी एक धारा जिसका घनत्व-अतिरेक \qty{10}{\percent} है, जो दिखाता है कि 1929 की धारा कितनी बड़ी और सघन रही होगी। (ग)~गति $\sqrt S$ के साथ बदलती है: $4.04/\sqrt 5 = \qty{1.81}{m.s^{-1}}$। धीमी धारा की प्रक्षुब्धता दुर्बल है और वह अपने मोटे कण गिरा देती है; वह कम सघन होती है, और धीमी होती है, और अपना भार मैदान पर एक \omterm{def:b1:clastic-sediments:turbidite}{क्रमित परत} के रूप में फैलाती है।
\end{solution}

\begin{solution}{exo:b1:clastic-sediments:12}
(क)~$0.7 \times \qty{5e10}{kg}/\qty{2000}{kg.m^{-3}} = \qty{1.75e7}{m^3}$ प्रति
वर्ष। (ख)~\qty{30}{m} ऊँचे और \qty{100}{km} चौड़े मुख पर फैलकर, वह मुख को
$\num{1.75e7}/(30 \times \num{1e5}) = \qty{5.8}{m}$ प्रति
वर्ष आगे बढ़ाता है। (ग)~$a = v\tan\alpha = 5.83 \times \tan\ang{0.5} = \qty{0.051}{m}$, लगभग
\qty{5}{cm} प्रति वर्ष, आगे बढ़ने से सौ गुना कम; \qty{30}{m} का क्रम $30/0.051 = 590$ वर्षों में बनता है।
\end{solution}

\begin{solution}{pb:b1:clastic-sediments:1}
\textbf{1.} इकाई~1: $(5.0 + 6.6 + 8.4)/3 = 6.67\,\varphi$, \qty{0.0098}{mm}
(\qty{9.8}{\micro m}); इकाई~2: $3.47\,\varphi$, \qty{0.090}{mm}; इकाई~3:
$1.80\,\varphi$, \qty{0.29}{mm}; इकाई~4: $0.73\,\varphi$, \qty{0.60}{mm}।

\textbf{2.} इकाई~1: $(8.4 - 5.0)/4 + (9.6 - 4.2)/6.6 = 0.850 + 0.818 =
1.67\,\varphi$; इकाई~2: $0.550 + 0.545 = 1.10$; इकाई~3: $0.275 + 0.273 = 0.55$;
इकाई~4: $0.725 + 0.712 = 1.44$।

\textbf{3.} इकाई~1: बारीक गाद, खराब छँटी; इकाई~2: बहुत बारीक बालू, खराब
छँटी; इकाई~3: मध्यम बालू, मध्यम रूप से अच्छी तरह छँटी; इकाई~4: मोटी बालू,
खराब छँटी।

\textbf{4.} इकाई~3। नदी-मुख पर धारा और लहरें बालू को बार-बार फिर से काम करती हैं और बारीकों को बहा ले जाती हैं, जबकि पूर्व-डेल्टा की कीचड़ निलंबन से उन सभी बारीक आकारों के मिश्रण के रूप में बैठती है जो वहाँ पहुँचते हैं।

\textbf{5.} आकार $8\,\varphi$, रास्ते के $1.4/1.8$ भाग पर, जो $\varphi_{50}$ से $\varphi_{84}$ तक जाता है, है, इसलिए $50 + (1.4/1.8) \times 34 = 76$ प्रतिशत
नमूना मोटा है और लगभग \qty{24}{\percent} मृत्तिका-आकार का है। ऐसे बारीक कण केवल लगभग स्थिर जल से बैठते हैं: स्थल लहरों और धाराओं की पहुँच के नीचे था।

\textbf{6.} एक धारा बजरी दोनों ले जाती है, जो बाढ़ों में उसके तल पर लुढ़कती है, और बालू; निक्षेप सबसे गहरे प्रवाह के बजरी-अवशेष को निचले प्रवाहों की बालू के साथ मिलाता है, इसलिए वह कई वर्गों को समेटता है।

\textbf{7.} माध्य \qty{0.0098}{mm} से \qty{0.090}{mm} और
\qty{0.29}{mm} तक बढ़ता है: क्रम ऊपर की ओर मोटा होता है।

\textbf{8.} इकाई~1, पटलित कीचड़: पूर्व-डेल्टा; इकाई~2, एकांतर गाद और बालू: निचला डेल्टा-मुख; इकाई~3, तिर्यक-परतदार बालू: ऊपरी डेल्टा-मुख की मुख-पट्टियाँ।

\textbf{9.} हर क्षेत्र उस पर आ गया जो पहले उसके सामने था: डेल्टा आगे बढ़ा, समुद्र की ओर, स्थल के ऊपर।

\textbf{10.} कम से कम \qty{21}{m}, इकाइयों 1 से 3 की मोटाई, क्योंकि उत्तराधिकार ने पानी को पूर्व-डेल्टा के तल से मुख-पट्टी तक भर दिया, सतह के निकट; और अधिक, क्योंकि तब से कीचड़ संहत हुई है (\cref{ch:b1:diagenesis-basins})।

\textbf{11.} डेल्टा-मैदान की एक वितरिका-धारा, जो डेल्टा के आगे बढ़ते समय मुख-पट्टी में कटी; उसकी बजरी उसका \omterm{def:b1:erosion-transport:load}{तल-भार} है।

\textbf{12.} नहीं: एक बिंदु-पट्टी एक अपरदन-सतह और एक अवशेष से शुरू होती है और बाढ़-मैदान की कीचड़ तक ऊपर की ओर बारीक होती है। डेल्टा-मुख कीचड़ से ऊपर की ओर मोटा होता है।

\textbf{13.} $\qty{16}{m}/\num{4000}\text{ वर्ष} = \qty{4.0}{mm}$ प्रति वर्ष।

\textbf{14.} $v = a/\tan\alpha = 0.0040/\tan\ang{0.5} = 0.0040/0.008727 =
\qty{0.458}{m}$ प्रति वर्ष।

\textbf{15.} $0.458 \times 4000 = \qty{1830}{m}$, लगभग \qty{1.8}{km}।

\textbf{16.} $\qty{21}{m}/\qty{4.0}{mm.yr^{-1}} = 5250$ वर्ष।

\textbf{17.} डेल्टा-मुख मुखों के निकट अधिकांश बालू पाता है, पूर्व-डेल्टा कहीं कम, और डेल्टा-मैदान बाढ़ों के अतिरिक्त थोड़ा; सक्रिय मुख भी एक जगह से दूसरी जगह खिसकते हैं, इसलिए एक बिंदु कुछ समय के लिए बहुत कुछ पा सकता है और फिर लगभग कुछ नहीं। दो तिथित परतों के बीच एक एकल दर इन परिवर्तनों पर एक माध्य है, और कीचड़ ने संहनन से अपनी मोटाई का कुछ भाग खो दिया है।

\textbf{18.} योग 1544 है, इसलिए अंकगणितीय माध्य \ang{128.7} है, एक दक्षिण-पूर्वी दिशा: निरर्थक, क्योंकि पाठन उत्तर के चारों ओर बिखरे हैं और 340 और 010 के निकट के मान ऐसे माध्य किए गए हैं मानो वे दूर हों।

\textbf{19.} $S = \sum\sin\alpha_i = 1.698$ और $C = \sum\cos\alpha_i =
11.248$।

\textbf{20.} $\bar R = \sqrt{1.698^2 + 11.248^2}/12 = 11.376/12 = 0.948$: पाठन एक दिशा के चारों ओर कसकर समूहित हैं।

\textbf{21.} नदियाँ लगभग उत्तर की ओर बहीं, और डेल्टा उसी ओर बना
(\cref{prop:b1:clastic-sediments:azimuths}); भूमि दक्षिण में थी।

\textbf{22.} क्योंकि अग्र-परतें उन परतों पर मापी गईं जो अब क्षैतिज नहीं हैं: उनकी नति-दिशाएँ झुकी परतों की, और झुकाव की, हैं। परतबंदी के \omterm{def:b1:geological-maps:attitude}{स्ट्राइक} के चारों ओर परतबंदी की \omterm{def:b1:geological-maps:attitude}{नति} से हर अग्र-परत घुमाने पर परतबंदी, और उसके साथ अग्र-परतें, अपने मूल \omterm{def:b1:geological-maps:attitude}{अभिविन्यास} में लौट आती हैं।

\textbf{23.} ज्वारीय धाराएँ अंदर-बाहर बहती हैं, और दो विपरीत दिशाओं में नत समुच्चय बनाती हैं, हेरिंगबोन \omterm{def:b1:clastic-sediments:cross-bedding}{तिर्यक-परत} की तरह; गुलाब-आरेख लगभग \ang{180} दूर दो पालि दिखाता। बारह पाठन एक पालि बनाते हैं, $\bar R$ 1 के निकट: एक नदी-धारा जो एक दिशा में बहती है।

\textbf{24.} $\bar\alpha = \operatorname{atan2}(1.698, 11.248) = \ang{8.6}$: इकाई~3 की माध्य \omterm{def:b1:clastic-sediments:cross-bedding}{पुरा-धारा} दिगंश \ang{9} है, उत्तर से थोड़ा पूर्व।
\end{solution}
